% Lösungen Einheit 2 von 3 – Binomische Formeln (zusammenbau v0.1)
\einheitenkopf{Einheit 2 von 3 · Binomische Formeln}

\erg{13}{a) ja: $(x + 9) \cdot (x + 9)$}
%% TODO Lösung der Erklärzeile 14g) fehlt in der Bank
\erg{14}{a) erste Formel; $a = x$, $b = 14$ \quad b) $x^2 + 8x + 16$ \quad c) $x^2 + 18x + 81$ \quad d) $x^2 + 2x + 1$ \quad e) $x^2 + 16x + 64$ \quad f) $x^2 + 22x + 121$ \quad g) \ldots \quad h) $x^2 - 18x + 81$ \quad i) $x^2 - 64$ \quad j) dritte Formel: $x^2 - 81$ \quad k) $9x^2 + 6x + 1$ \quad l) $x^2 + 6xy + 9y^2$ \quad m) $2 \cdot (x^2 + 8x + 16) = 2x^2 + 16x + 32$ \quad n) $-x^2 - 18x - 81$ \quad o) $x^2 + 2x + 1 - 6 = x^2 + 2x - 5$ \quad p) $x^2 + 8x + 16 - 7 = x^2 + 8x + 9$; die Gleichungen stimmen überein.}
\erg{15}{a) $(40 + 1)^2 = 1\,600 + 80 + 1 = 1\,681$}
\erg{16}{a) Fehler: das Mittelglied fehlt. Richtig: $y^2 + 18y + 81$. \quad b) Für $x = 2$: links $6^2 = 36$, rechts $4 + 16 = 20$. Beim Quadrieren einer Summe entsteht das Mittelglied $8x$, das Tim weglässt.}
